\documentclass[a4paper,11pt,oneside]{memoir}

% Castellano
\usepackage[spanish,es-tabla]{babel}
\selectlanguage{spanish}
\usepackage[utf8]{inputenc}
\usepackage{placeins}
\usepackage{float}

\RequirePackage{booktabs}
\RequirePackage[table]{xcolor}
\RequirePackage{xtab}
\RequirePackage{multirow}

% Links
\usepackage[colorlinks]{hyperref}
\hypersetup{
	allcolors = {red}
}

% Ecuaciones
\usepackage{amsmath}

% Rutas de fichero / paquete
\newcommand{\ruta}[1]{{\sffamily #1}}

% Párrafos
\nonzeroparskip


% Imagenes

%\usepackage{graphicx}
%\newcommand{\imagen}[2]{
%	\begin{figure}[!h]
%		\centering
%		\includegraphics[width=0.9\textwidth]{#1}
%		\caption{#2}\label{fig:#1}
%	\end{figure}
%	\FloatBarrier
%}

\newcommand{\imagenflotante}[2]{
	\begin{figure}%[!h]
		\centering
		\includegraphics[width=0.9\textwidth]{#1}
		\caption{#2}\label{fig:#1}
	\end{figure}
}

% Links
\usepackage[colorlinks]{hyperref}
\hypersetup{
	colorlinks,
	linkcolor={green!40!black},
	citecolor={blue!50!black},
	urlcolor={blue!80!black}
}


% El comando \figura nos permite insertar figuras comodamente, y utilizando
% siempre el mismo formato. Los parametros son:
% 1 -> Porcentaje del ancho de página que ocupará la figura (de 0 a 1)
% 2 --> Fichero de la imagen
% 3 --> Texto a pie de imagen
% 4 --> Etiqueta (label) para referencias
% 5 --> Opciones que queramos pasarle al \includegraphics
% 6 --> Opciones de posicionamiento a pasarle a \begin{figure}
\newcommand{\figuraConPosicion}[6]{%
  \setlength{\anchoFloat}{#1\textwidth}%
  \addtolength{\anchoFloat}{-4\fboxsep}%
  \setlength{\anchoFigura}{\anchoFloat}%
  \begin{figure}[#6]
    \begin{center}%
      \Ovalbox{%
        \begin{minipage}{\anchoFloat}%
          \begin{center}%
            \includegraphics[width=\anchoFigura,#5]{#2}%
            \caption{#3}%
            \label{#4}%
          \end{center}%
        \end{minipage}
      }%
    \end{center}%
  \end{figure}%
}

% Imagenes
\usepackage{graphicx}
\newcommand{\imagen}[2]{
	\begin{figure}[!h]
		\centering
		\includegraphics[width=0.9\textwidth]{#1}
		\caption{#2}\label{fig:#1}
	\end{figure}
	\FloatBarrier
}

%\newcommand{\imagenflotante}[2]{
%	\begin{figure}%[!h]
%		\centering
%		\includegraphics[width=0.9\textwidth]{#1}
%		\caption{#2}\label{fig:#1}
%	\end{figure}
%}

%
% Comando para incluir imágenes en formato apaisado (sin marco).
\newcommand{\figuraApaisadaSinMarco}[5]{%
  \begin{figure}%
    \begin{center}%
    \includegraphics[angle=90,height=#1\textheight,#5]{#2}%
    \caption{#3}%
    \label{#4}%
    \end{center}%
  \end{figure}%
}
% Para las tablas
\newcommand{\otoprule}{\midrule [\heavyrulewidth]}
%
% Nuevo comando para tablas pequeñas (menos de una página).
\newcommand{\tablaSmall}[5]{%
 \begin{table}
  \begin{center}
   \rowcolors {2}{gray!35}{}
   \begin{tabular}{#2}
    \toprule
    #4
    \otoprule
    #5
    \bottomrule
   \end{tabular}
   \caption{#1}
   \label{tabla:#3}
  \end{center}
 \end{table}
}

%
% Nuevo comando para tablas pequeñas (menos de una página).
\newcommand{\tablaSmallSinColores}[5]{%
 \begin{table}[H]
  \begin{center}
   \begin{tabular}{#2}
    \toprule
    #4
    \otoprule
    #5
    \bottomrule
   \end{tabular}
   \caption{#1}
   \label{tabla:#3}
  \end{center}
 \end{table}
}

\newcommand{\tablaApaisadaSmall}[5]{%
\begin{landscape}
  \begin{table}
   \begin{center}
    \rowcolors {2}{gray!35}{}
    \begin{tabular}{#2}
     \toprule
     #4
     \otoprule
     #5
     \bottomrule
    \end{tabular}
    \caption{#1}
    \label{tabla:#3}
   \end{center}
  \end{table}
\end{landscape}
}

%
% Nuevo comando para tablas grandes con cabecera y filas alternas coloreadas en gris.
\newcommand{\tabla}[6]{%
  \begin{center}
    \tablefirsthead{
      \toprule
      #5
      \otoprule
    }
    \tablehead{
      \multicolumn{#3}{l}{\small\sl continúa desde la página anterior}\\
      \toprule
      #5
      \otoprule
    }
    \tabletail{
      \hline
      \multicolumn{#3}{r}{\small\sl continúa en la página siguiente}\\
    }
    \tablelasttail{
      \hline
    }
    \bottomcaption{#1}
    \rowcolors {2}{gray!35}{}
    \begin{xtabular}{#2}
      #6
      \bottomrule
    \end{xtabular}
    \label{tabla:#4}
  \end{center}
}

%
% Nuevo comando para tablas grandes con cabecera.
\newcommand{\tablaSinColores}[6]{%
  \begin{center}
    \tablefirsthead{
      \toprule
      #5
      \otoprule
    }
    \tablehead{
      \multicolumn{#3}{l}{\small\sl continúa desde la página anterior}\\
      \toprule
      #5
      \otoprule
    }
    \tabletail{
      \hline
      \multicolumn{#3}{r}{\small\sl continúa en la página siguiente}\\
    }
    \tablelasttail{
      \hline
    }
    \bottomcaption{#1}
    \begin{xtabular}{#2}
      #6
      \bottomrule
    \end{xtabular}
    \label{tabla:#4}
  \end{center}
}

%
% Nuevo comando para tablas grandes sin cabecera.
\newcommand{\tablaSinCabecera}[5]{%
  \begin{center}
    \tablefirsthead{
      \toprule
    }
    \tablehead{
      \multicolumn{#3}{l}{\small\sl continúa desde la página anterior}\\
      \hline
    }
    \tabletail{
      \hline
      \multicolumn{#3}{r}{\small\sl continúa en la página siguiente}\\
    }
    \tablelasttail{
      \hline
    }
    \bottomcaption{#1}
  \begin{xtabular}{#2}
    #5
   \bottomrule
  \end{xtabular}
  \label{tabla:#4}
  \end{center}
}



\definecolor{cgoLight}{HTML}{EEEEEE}
\definecolor{cgoExtralight}{HTML}{FFFFFF}

%
% Nuevo comando para tablas grandes sin cabecera.
\newcommand{\tablaSinCabeceraConBandas}[5]{%
  \begin{center}
    \tablefirsthead{
      \toprule
    }
    \tablehead{
      \multicolumn{#3}{l}{\small\sl continúa desde la página anterior}\\
      \hline
    }
    \tabletail{
      \hline
      \multicolumn{#3}{r}{\small\sl continúa en la página siguiente}\\
    }
    \tablelasttail{
      \hline
    }
    \bottomcaption{#1}
    \rowcolors[]{1}{cgoExtralight}{cgoLight}

  \begin{xtabular}{#2}
    #5
   \bottomrule
  \end{xtabular}
  \label{tabla:#4}
  \end{center}
}


















\graphicspath{ {./img/} }

% Capítulos
\chapterstyle{bianchi}
\newcommand{\capitulo}[2]{
	\setcounter{chapter}{#1}
	\setcounter{section}{0}
	\chapter*{#2}
	\addcontentsline{toc}{chapter}{#2}
	\markboth{#2}{#2}
}

% Apéndices
\renewcommand{\appendixname}{Apéndice}
\renewcommand*\cftappendixname{\appendixname}

\newcommand{\apendice}[1]{
	%\renewcommand{\thechapter}{A}
	\chapter{#1}
}

\renewcommand*\cftappendixname{\appendixname\ }

% Formato de portada
\makeatletter
\usepackage{xcolor}
\newcommand{\tutor}[1]{\def\@tutor{#1}}
\newcommand{\course}[1]{\def\@course{#1}}
\definecolor{cpardoBox}{HTML}{E6E6FF}
\def\maketitle{
  \null
  \thispagestyle{empty}
  % Cabecera ----------------
\noindent\includegraphics[width=\textwidth]{cabecera}\vspace{1cm}%
  \vfill
  % Título proyecto y escudo informática ----------------
  \colorbox{cpardoBox}{%
    \begin{minipage}{.8\textwidth}
      \vspace{.5cm}\Large
      \begin{center}
      \textbf{TFG del Grado en Ingeniería Informática}\vspace{.6cm}\\
      \textbf{\LARGE\@title{}}
      \end{center}
      \vspace{.2cm}
    \end{minipage}

  }%
  \hfill\begin{minipage}{.20\textwidth}
    \includegraphics[width=\textwidth]{escudoInfor}
  \end{minipage}
  \vfill
  % Datos de alumno, curso y tutores ------------------
  \begin{center}%
  {%
    \noindent\LARGE
    Presentado por \@author{}\\ 
    en Universidad de Burgos --- \@date{}\\
    Tutor: \@tutor{}\\
  }%
  \end{center}%
  \null
  \cleardoublepage
  }
\makeatother

\newcommand{\nombre}{Adrian Aguado} %%% cambio de comando

% Datos de portada
\title{{\Huge barterAPP}\\[0.5cm]Control to your business time..}
\author{\nombre}
\tutor{Luisz R. Izquierdo}
\date{\today}

\begin{document}

\maketitle


\newpage\null\thispagestyle{empty}\newpage


%%%%%%%%%%%%%%%%%%%%%%%%%%%%%%%%%%%%%%%%%%%%%%%%%%%%%%%%%%%%%%%%%%%%%%%%%%%%%%%%%%%%%%%%
\thispagestyle{empty}


\noindent\includegraphics[width=\textwidth]{cabecera}\vspace{1cm}

\noindent D. Luis R. izquierdo, profesor del departamento de Ingeniería Civil área de organización de empresas

\noindent Expone:

\noindent Que el alumno D. \nombre, con DNI 78759314T, ha realizado el Trabajo final de Grado en Ingeniería Informática titulado  barterAPP 

\noindent Y que dicho trabajo ha sido realizado por el alumno bajo la dirección del que suscribe, en virtud de lo cual se autoriza su presentación y defensa.

\begin{center} %\large
En Burgos, {\large \today}
\end{center}

\vfill\vfill\vfill

% Author and supervisor
%\begin{minipage}{0.45\textwidth}
%\begin{flushleft} %\large
%Vº. Bº. del Tutor:\\[2cm]
%D. nombre tutor
%\end{flushleft}
%\end{minipage}
%\hfill
%\begin{minipage}{0.45\textwidth}
%\begin{flushleft} %\large
%Vº. Bº. del co-tutor:\\[2cm]
%D. nombre co-tutor
%\end{flushleft}
%\end{minipage}
%\hfill
%
%\vfill

% para casos con solo un tutor comentar lo anterior
% y descomentar lo siguiente
Vº. Bº. del Tutor:\\[2cm]
Luis R. Izquierdo


\newpage\null\thispagestyle{empty}\newpage




\frontmatter

% Abstract en castellano
\renewcommand*\abstractname{Resumen}
\begin{abstract}
La aplicacion barterAPP mola un montón por que
\end{abstract}

\renewcommand*\abstractname{Descriptores}
\begin{abstract}
WebApp, Angular2, Nodejs, Express
\end{abstract}

\clearpage

% Abstract en inglés
\renewcommand*\abstractname{Abstract}
\begin{abstract}
BarterApp is really great because is super
\end{abstract}

\renewcommand*\abstractname{Keywords}
\begin{abstract}
WebApp, Angular2, Nodejs, Express
\end{abstract}

\clearpage

% Indices
\tableofcontents

\clearpage

\listoffigures

\clearpage

\listoftables
\clearpage

\mainmatter
\capitulo{1}{Introducción}


Hoy en día existen multitud de situaciones en las que se requiere un reparto de turnos o guardias entre un grupo de personas (por ejemplo personal sanitario en servicios hospitalarios, servicios de emergencias extrahospitalarias, bomberos, protección civil y muchos). En muchas de éstas ocasiones, el reparto de turnos planificado podría mejorarse fácilmente si se llevaran a cabo cambios bilaterales que a veces no se llegan a producir porque las partes implicadas desconocen que existe esa oportunidad de mejora, o porque intuyen que los costes de encontrar esa mejora y llegar a  hacerla posible son excesivamente elevados. Otras veces los cambios sí que pueden llegar a producirse pero tras un largo período de negociación.

Actualmente, tras un proceso de investigación inicial, he determinado que no existe una herramienta sencilla que presente una solución a este problema. Es más en servicios como los sanitarios el problema no esta ni si quiera digitalizado, es más se suele realizar simple y llanamente en papel. Por una parte esta no es una mala técnica pero la implantación de una herramienta digital accesible por cualquiera puede facilitar las cosas enormemente a nivel de tiempo, espacio y entendimiento. 

El mundo de las aplicaciones, tanto web como móvil, resulta extremadamente competitivo hoy en día y posicionarse en el mercado es realmente complejo. Por una parte esta la tarea de adaptarse a las nueva tecnologías tanto web, prácticamente cada muy pocos meses cambian las tecnologías o los frameworks con los que realizar las webs; como también móvil, versiones que se van actualizando de todos los sistemas operativos. 

\section{Motivacion}\label{Motivacion}

Explicar por que he elegido este trabajo y no otro es también describir cuáles son mis objetivos personales de cara a un futuro cada vez más cercano.  Me gusta el desarrollo web y me apasiona la posibilidad de mejorar el mundo a través de la tecnología, si bien no tengo experiencia dentro del desarrollo web si que me gustaría encaminarme hacia ello en mis próximos años. 

\section{Estructura de la memoria}\label{estructura-de-la-memoria}

\begin{itemize}
\tightlist
\item
  \textbf{Introducción:} Aquí se explica la motivación que me ha llevado a escoger este trabajo y no otro, la descripción del problema propuesto y una breve introducción general a la solución que se ha pretendido dar. Así mismo también una estructura de toda la memoria 
\item
  \textbf{Objetivos del proyecto:} tanto a nivel académico, personal como los objetivos alcanzados con el proyecto. También se comentará lo que le espera en un futuro al proyecto, ya que me gustaría seguir con él.
\item
  \textbf{Conceptos teóricos:} breve explicación de los conceptos
  teóricos que he debido adquirir previamente para la realización del proyecto. Necesarios para comprender mejor cómo funciona el proyecto.
\item
  \textbf{Técnicas y herramientas:} software principal, metodologías empleadas durante el proyecto. frameworks empleados y herramientas.
\item
  \textbf{Aspectos relevantes del desarrollo:} detalles sobre el la realización del proyecto.
\item
  \textbf{Trabajos relacionados:} otros aspectos del desarrollo del proyecto.
\item
  \textbf{Conclusiones y líneas de trabajo futuras:} conclusiones
  obtenidas tras la realización de este trabajo fin de grado y hacia donde va en el futuro.
\end{itemize}

Junto a la memoria se proporcionan los siguientes anexos:

\begin{itemize}
\tightlist
\item
  \textbf{Plan del proyecto software:} estudio de viabilidad y planificación del proyecto paso a paso.
\item
  \textbf{Especificación de requisitos del software:} se describe la
  fase de análisis; los objetivos generales, el catálogo de requisitos
  del sistema y la especificación de requisitos necesarios para su correcto funcionamiento.
\item
  \textbf{Especificación de diseño:} se describe la fase de diseño de la aplicación tanto front-end como back-end y el paso de la interfaz web a la móvil.
\item
  \textbf{Manual del programador:} recoge los aspectos más relevantes
  relacionados con el código fuente y su correcto manejo por parte de un programador.
  \textbf{Manual de usuario:} guía de usuario para el uso de la aplicación.
\end{itemize}


\capitulo{2}{Objetivos del proyecto}

En este apartado se van a detallar los diferentes objetivos que se buscaban, a diferentes niveles, con la realización del proyecto.

\section{Objetivos generales}\label{objetivos-generales}
Comenzaremos con los objetivos generales del proyecto. 
\begin{itemize}
\tightlist
\item
  Desarrollar una  \emph{aplicación web} que permita la gestión de turnos de una manera sencilla e intuitiva. 
\item
  Desarrollar una aplicación móvil híbrida que permita acceder desde cualquier plataforma a esa aplicación para dar un servicio multiplataforma. 
\item
  Permitir tener una herramienta que resuelva un problema real.
\end{itemize}

\section{Objetivos técnicos}\label{objetivos-tecnicos}
A continuación explicaré los principales objetivos técnicos que se pretendían. 
\begin{itemize}
\tightlist
\item
  Toma de contacto con el mundo web por parte del alumno, a su juicio el dónde . 
 \item
  Aprendizaje de la implementación de una aplicación web con  \emph{frameworks} web, como puede ser el caso de Angular o Ionic. 
\item
  Breve toma de contacto con el mundo móvil (\emph{PhoneGapp, Cordova}).
\item
 Aplicar la metología Scrum.
 \item
 Utilizar ZenHub como Herramienta de gestión de proyectos
\end{itemize}

\section{Objetivos personales}\label{objetivos-personales}
Los objetivos personales que he perseguido durante todo el desarrollo han sido los siguientes:
\begin{itemize}
\tightlist
\item
 El diseño e implementación de una aplicación web, realizada con \emph{Node.js}y \emph{Express} para la parte del
servidor, y \emph{Angular} y \emph{Bootstrap} para la parte del cliente. 
\item
  Conocer y manejar bases de datos NoSQL, como \emph{MongoDB}.
\item
  Mejorar mi inquietud y curiosidad sobre el mundo web. 
\item
  Averiguar y aprender el mayor número de cosas sobre las tecnologías web del mercado actual. 
\item
  Aprender a realizar un proyecto de gran envergadura, y sobre todo, a gestionarlo.
\end{itemize}

\section{Objetivos alcanzados}\label{objetivos-alcanzados}
Los objetivos que finalmente se han alcanzado han sido: 
\begin{itemize}
\tightlist
\item
  Aplicación web funcionando. 
\item
  Desarrollar una \emph{aplicación web} . 
\item
  Hola hola.
\end{itemize}

\section{Objetivos futuros}\label{objetivos-futuros}
Esta sección es una breve enumeración de lo que se pretende detallar en la última parte de la memoria.

\begin{itemize}
\tightlist
\item
  Lograr una mejora sustanciar de la usabilidad de la aplicación en su versión móvil. 
\item
  Realizar encuesta para mejorar la aplicación base en función de lo que lo usuarios demanden . 
\item
  Implementar un chat que permita poner en contacto directamente a los usuarios.
\item
  Posicionar la herramienta dentro del mercado.
\end{itemize}
\capitulo{3}{Conceptos teóricos}

La parte del proyecto más ardua ha sido, sin lugar a dudas, la de aprender a utilizar \emph{frameworks}  web, para lo cual me he nutrido de diversos cursos online, tanto gratuitos como de pago. A continuación, en este apartado voy a relatar todas las tecnologías empleadas y a justificar el por que he empleado unas y no otras. Las técnicas y tecnologías como tal se verán en el siguiente apartado pero todas ellas presentan unos conceptos teóricos comunes que es necesario conocer. 


\section{Introducción}\label{teorico-introduccion}
El desarrollo de aplicaciones informáticas evoluciona continuamente para adaptarse a las tecnologías de la información y las comunicaciones (TIC). El auge de Internet y de la web ha influido notablemente en el desarrollo de software durante los últimos años. Hoy en día la interfaz de los sistemas de información se implementa utilizando tecnologías web que ofrecen numerosas ventajas tales como el uso de una interfaz uniforme y la mejora del mantenimiento del sistema. Sin embargo, la existencia de numerosos estándares y los intereses de los fabricantes de tecnologías web dificultan el desarrollo de este tipo de aplicaciones. 

En los principios de la informática las relaciones en entre los usuarios y los programas de los que hacían uso era muy diferente a lo que tenemos ahora. Ahora se potencia el diseño basado en usuario, antes es potencia más que el programa estuviera implementado de manera idónea antes que la experiencia que el usuario tenía al hacer uso del mismo. Obviamente ahora también se impulsa esa manera de desarrollar código, lo que se conoce como \emph{clean code} \cite{https://g.co/kgs/VLR4mn}, pero en un mundo en el que la competitividad es tan alta prima sobre todo la experiencia del usuario final, que a fin y al cabo va a ser el que interactúa con tu producto ya sea web o móvil. 

\imagen{uxschema}{Esquema \emph{UI, UX IxD}. Fuente  \url{dealfuel.com}.}



El mundo de las aplicaciones web es un mundo de constante evolución y actualización, donde aparecen nuevas tecnologías que ofrecen tanto mejoras visuales que mejoran la experiencia del usuario como de rendimiento. Este tipo de tecnologías abren un abanico inmenso de posibilidades, y hacen pensar a los analistas y programadores de aplicaciones web que puede haber otras soluciones que mejoren su aplicación, pero que no estaban disponibles cuando definieron la arquitectura y el diseño de sus aplicaciones. Es un mundo que avanza tan rápido que en muchas ocasiones es imposible estar al día de todos los  \emph{frameworks} o tecnologías que van surgiendo, por eso es mejor focalizar los esfuerzos en alguna en concreto e intentar aprenderla de la mejor manera posible. 




\section{Concepto general del proyecto}\label{teorico-general}
Hace algún tiempo para realizar una web existía una gran barrera a la hora de entender el concepto de cliente y servidor. Por un lado estaba la parte del cliente, la cual se realizaba en lenguajes puros que todos conocemos como \emph{HTML} y \emph{CSS} para las hojas de estilo. Por otro lado estaba el servidor lo cuál significa cambiar totalmente de lenguaje, lo que suponía un salto para un programador web que debía conocer ambos lados para crear webs seguras y robustas, la parte del cliente es la que interactúa con el usuario, la que nosotros vemos, y la parte del servidor es la parte que conecta con la base de datos, en caso de que sea necesario.

Todo cambio cuando se popularizó Javascript para la realización webapps, la posibilidad de crear web con un mismo lenguaje en todas las partes resulta muy atractiva tanto para el programador como para la lógica del propio programa o web ya que se disminuye el número de errores o se consiguen localizar de manera más sencilla al no tener que estar cambiando de lenguaje. En este punto es donde surge el stack MEAN:

\imagen{meanjs}{Esquema \emph{MEAN}. Fuente: \url{kambrica.com}.} 



La presencia de \emph{Javascript} para el desarrollo de software se está haciendo un hueco cada vez mayor en el mercado. Aquel humilde lenguaje que empezó en los años '90 como una vía sencilla de validar formularios, se ha convertido en parte fundamental del desarrollo de todo tipo de aplicaciones: web, móviles, bases de datos e, incluso, administración de sistemas. Esta proliferación ha llevado a  \emph{Javascript} a todas las capas de desarrollo, empezando por el lado cliente en sus inicios (el navegador), pero yendo también al servidor y a la capa de almacenamiento. En cualquiera de esos puntos podemos encontrar  \emph{Javascript} listo para ser utilizado.







\section{Análisis del sistema: MVC}\label{analisis-sistema-mvc}

El \emph{modelo-vista-controlador} un patrón de arquitectura de software, que separa los datos y la lógica de negocio de una aplicación de la interfaz de usuario y el módulo encargado de gestionar los eventos y las comunicaciones. Para ello propone la construcción de tres componentes distintos que son el modelo, la vista y el controlador, es decir, por un lado define componentes para la representación de la información, y por otro lado para la interacción con el usuario.

\imagen{mvcschema2}{Dibujo \emph{MVC}\cite{mvc}. Fuente: \url{artima.com}.}




\section{Arquitectura}\label{arquitectura}

Como se ha comentado en secciones anteriores la aplicación tiene dos sistemas bien diferenciados: servidor y cliente. En esta sección se expondrán las arquitecturas de cada una de ellas así como de la aplicación competente.

\capitulo{4}{Técnicas y herramientas}

Esta parte de la memoria tiene como objetivo presentar las tecnologías y las herramientas de desarrollo que se han utilizado para llevar a cabo el proyecto. Se han estudiado diferentes alternativas de metodologías, herramientas, bibliotecas  y se pretende aquí realizar un resumen de los aspectos más destacados de cada alternativa, incluyendo comparativas entre las distintas opciones y una justificación de las elecciones realizadas. 

No se pretende que este apartado se convierta en un capítulo de un libro dedicado a cada una de las alternativas detalladamente, sino comentar los aspectos más destacados de cada opción.

He de decir que durante el desarrollo del proyecto tanto las herramientas cómo las técnicas a utilizar han ido mutando debido a que yo mismo iba descubriendo nuevas técnicas o nuevas herramientas que mejoraban mi manera de hacer el código o me facilitaban mi manera de trabajar. No ha sido así con la tecnología ya que esto hubiera supuesto más problemas que ventajas. El estar cambiando de herramienta puede que haya sido uno de los mayores errores del proyecto dado que siempre es mejor focalizarse en una sola antes de intentar abarcar demasiado. 


\section{Metodologías}\label{metodologias}
En este apartado se describen las metodologías utilizadas para el desarrollo del sistema.

\subsection{Scrum}\label{scrum}
Scrum es un sistema de desarrollo de software que está dentro de las metologías ágiles. Principalmente, se basa en la creación y asignación de tareas, normalmente escribiéndolas en un postit de determinado color. Esta tarea entrara en un sistema iterativo que controlará los cambios de estado de la tarea, hasta que esta resulta completada y se descarta del panel de scrum.

\subsection{MVC: Model-View-Controller}\label{mvc}
Ya nombrado anteriormente, Ees un patrón de diseño software utilizada para implementar sistemas donde se requiere el uso de interfaces de usuario. Surge de la necesidad de crear software más robusto con un ciclo de vida más adecuado, donde se potencie la facilidad de mantenimiento, reutilización del código y la separación de conceptos. Es decir surge de la idea de separar el modelo por un lado, la vista por otro y por último el controlador.


\imagen{mvcschema}{Esquema \emph{MVC}. Elaboración propia.}

\subsection{Patrones de diseño: singleton}\label{Singleton}


El patrón de diseño \emph{singleton} \cite{wikisingleton} está diseñado para restringir la creación de objetos pertenecientes a una clase o el valor de un tipo a un único objeto. Su intención consiste en garantizar que una clase sólo tenga una instancia y proporcionar un punto de acceso global a ella. Es decir se crean objetos, declarados como tipo privados, una sola vez para se usados en muchos sitios diferentes.

\section{Herramientas}\label{herramientas}
En este apartado se describen las herramientas utilizadas para el desarrollo del sistema.

\subsection{Atom}\label{herramientas_atom}
\hyperlink{https://atom.io/}{Atom} es un editor de texto moderno de código abierto desarrollado por  \hyperlink{http://brackets.io/}{GitHub}  que se ha escogido en este proyecto para desarrollar la aplicación debido a que esta desarrollado utilizando tecnologías web, luego ha sido creado por y para la web. Es posible ampliar sus funcionalidades a través de plugins desarrollados con \emph{Node.js} que puede ser instalados de una forma sencilla a través del gestor de paquetes internos con el que cuenta, así como diferentes tipos de temas. Esta posibilidad hace que se convierta en un editor de texto muy personal, ya que es el mismo desarrollador el que elige las características que desea tener sin afectar mínimamente al rendimiento.

\imagenflotante{logo-atom}{Logo \emph{Atom}. Fuente: \url{atom.io}.}

\begin{itemize}
	\item \subsubsection{Alternativas estudiadas: otros}
	\begin{itemize}
		\item \hyperlink{http://brackets.io/}{Brackets}
		\item \hyperlink{http://eclipse.org}{Eclipse}
	\end{itemize}
	Tal y como he nombrado en la introducción la necesidad o la curiosidad por mejorar mi trabajo hicieron que encontrara nuevas herramientas durante el desarrollo del proyecto, es el caso, por ejemplo, del entorno de desarrollo. Al comienzo empecé con Brackets como IDE, ya que era el entorno que utilizaba el profesor de los cursos que me recomendó  mi tutor. Después y atraído por el uso a diario de Eclipse como herramienta principal de trabajo durante mi período de prácticas con la universidad decidí que podía ser una buena idea el emplearlo también para realizar la API, /underline{craso error} dado que Eclipse no fue creado para el mundo del desarrollo web por lo que me vi con que cada vez me encontraba con menos soporte para lenguajes, sin ir más lejos eclipse no presenta, actualmente, soporte para  \emph{typescript} el lenguaje principal de desarrollo de  \emph{angular 2}. Por fin di con el IDE perfecto para desarrollar entornos web que es el descrito en el párrafo anterior.
\end{itemize}

\subsubsection{Detalles: WebStorm}\label{detalle_webstorm}
Otro \emph{IDE} que recomiendo y que ofrece una grandes prestaciones es \hyperlink{https://www.jetbrains.com/webstorm/}{Webstorm}. Según algunos artículos \cite{articulo1} es una de las mejores herramientas para Javascript del mercado, esta enfocado totalmente a este lenguaje web y sigue la misma filosofía de paquetes que Atom. 


\subsection{Robomongo}\label{herramientas_robomongo}
\hyperlink{http://http//robomongo.org/}{Robomongo} es una interfaz para \emph{MongoDB} que nos permite conectarnos al servidor de base de datos de forma sencilla ya que nada más arrancarlo podemos crear una nueva conexión. Resulta también muy sencillo e intuitivo de utilizar.
\imagenflotante{logo-robomongo}{Logo \emph{Robomongo}. Fuente: \url{robomongo.org}.}



\subsection{Firefox for developers}\label{herramientas_firefox}
Para el navegador he elegido Firefox en su versión para desarrolladores: \hyperlink{https://www.mozilla.org/es-ES/firefox/developer}{Firefox Developer Edition} . El cuál considero que posee todas las herramientas esenciales para testear y probar una aplicación web.
\imagenflotante{logo-firefox}{Logo \emph{FF Developers}. Fuente: \url{firefox.com/developer}.}


\section{Tecnologías}\label{tecnologias}
 En este apartado se describen las tecnologías utilizadas para el desarrollo del sistema. Las siguientes líneas servirán al lector para conocer algunas de las tecnologías más modernas y poco conocidas que se han usado en el proyecto. Se obvian el nombrar directamente lenguajes web como pueden ser \emph{Javascript}, \emph{HTML} o \emph{CSS}. Existe mucha información accesible en la web sobre ellas y se hará más incapié en las tecnologías principales que usan propiamente éstos lenguajes que no son tan comunes.
 
 \subsection{Front-end: Angular 2+}\label{tecnologias_angular}
 \hyperlink{https://angular.io}{Angular} es un framework cliente MVC \emph{Javascript} de código abierto creado en sus inicios por Google, permite crear  \emph{Single-Page Applications} (\hyperlink{https://es.wikipedia.org/wiki/Single-page_application}{SPA})  cuya principal característica es la de dar al usuario la impresión de que todo sucede en la misma página, sin hacer recargas de la misma. Es decir, trata de emular a las aplicaciones de escritorio, lo que se trata es de intercambiar las vista y no de recargar la página. En la actualidad cuenta con una amplia comunidad de desarrolladores  que dan soporte al framework.  Algunas características de Angular son:
 
\begin{itemize}
\tightlist
\item
  El sistema de databinding es muy completo y potente.
\item
  Angular es una solución completa que incluye prácticamente todos los aspectos que puedes necesitar para crear una aplicación cliente en \emph{Javascript} . 
\item
  La comunidad de desarrolladores crece cada día y la popularidad de Angular es un hecho.
  \item
  Está realizado para ser fácilmente testeable.
\end{itemize}
 
 \imagenflotante{logo-angular}{Logo \emph{Angular}. Fuente: \url{angular.io}.}
 

\subsubsection{Detalles: Angular Arquitectura}\label{detalle_angulararquitectura}
Un app de Angular esta basada en componentes. En su día resultó una revolución en el desarrollo web ya que lo que se persigue es que cada parte de la aplicación posea un componente de manera que se consigue hacer aplicaciones web reutilizables de algún modo.

\imagen{angularschema}{Esquema componentes. Fuente:  \url{rldona.gitbooks.io}.}


\begin{itemize}
\tightlist
\item
Vistas: es la parte visible por el usuario. Suelen estar parametrizadas, y todas
las vistan tienen asociado un controlador.
\item
 Controladores: Sirven los datos y funcionalidades a las vistas asociadas a él. Suelen ser estos los que utilizan los servicios para obtener los datos.
\item
 Directivas: Ofrecen elementos nuevos o nuevos comportamientos en elementos ya existentes dentro de las propias vistas.
  \item
  Servicios: son los encargados de proveer los datos. Lo más normal es que estos datos provengan de una API externa.
\end{itemize}


\subsubsection{Detalles: Typescript}\label{detalle_typescript}
Angular 4 esta basado en  \hyperlink{https://es.wikipedia.org/wiki/TypeScript}{Typescript} , es un lenguaje de código abierto que resulta ser un superset de Javascript. Es fuertemente tipado y orientado a objetos basado en clases.


\subsubsection{Detalles: Angular CLI}\label{detalle_angularcli}
Es un intérprete de línea de comandos de Angular que te facilitará el inicio y desarrollo de proyectos, ocupándose de la creación del esqueleto de la mayoría de los componentes de una aplicación Angular. Es interesante utilizarlo dado que ahorra mucho trabajo al desarrollador además de evitar despistes innecesarios. Por poner un ejemplo cada vez que un componente es creado en Angular es necesario modificar el componente principal manualmente para poder usarlo, con Angular CLI tan solo es necesario introducir un comando para crear el nuevo componente y Angular añade todas las dependencias por ti. Se podría pensar que se trata de una heramienta de terceros pero no es asi, la interfaz de comandos es proporcionada directamente por el equipo de Angular. Se darán consejos de como usarlo en los anexos.

Por otro lado el uso de Angular CLI permite realizar en un solo comando el preparado de la aplicación para lanzarla al servidor, lo cual resulta también muy interesante de cara a una parte más seria de la aplicación como puede ser por ejemplo lanzarla al mercado.

El aspecto de una aplicación hola mundo creada por Angular CLI contendría los siguientes elementos:

\imagenflotante{angularclihelloworld}{Vista \emph{App hola Mundo}. Elaboración propia.}

Más información en el Anexo: Manual del programador. 

\subsubsection{Detalles: Calendario Angular}\label{detalle_calendario}
El calendario elegido ha sido encontrado en github y realizado como un componente de Angular2+, es heredado de un calendario para AngularJS que ya tuvo mucho éxito y posteriormente realizado para versiones superiores. Es un calendario flexible y pensado para muchos usos por lo que no es completamente \emph{responsive} tal y cómo se puede ver en la documentación \cite{calendariodocs}  que está disponible y nos explica de manera detallada como incluirlo en nuestro proyecto Angular.

\begin{itemize}
	\item \subsubsection{Alternativas estudiadas: otros calendarios} 
	\begin{itemize}
		\item \hyperlink{https://fullcalendar.io/}{Full Calendar (Javascript)}
		\item \hyperlink{https://github.com/nekken/ng2-fullcalendar}{El componente de FullCalendar para Angular2}
		\item \hyperlink{https://doc.daypilot.org/calendar/angular-2/}{Day pilot }
		\item \hyperlink{http://angular-ui.github.io/ui-calendar/}{Calendario UI} interesante pero no disponible en versiones de angular 2+.
	\end{itemize}
	\end{itemize}

\subsubsection{Detalles: Módulos adicionales empleados con Angular2+}\label{detalle_angularmodulos}
Algunos de los módulos adicionales a destacar empleados en la parte del cliente han sido:

\begin{itemize}
\tightlist
\item
1-
\item
2-
\item
3-
  \item
4-
\end{itemize}

\begin{itemize}
	\item	\subsubsection{Alternativas estudiadas}
	
	Algunas de las alternativas estudiadas a Angular 2+ han sido las siguientes:
	\begin{description}
\item[ \hyperlink{backbonejs.org/}{BackBone JS} ] Es un framework web que permite construir aplicaciones usando Javascript siguiendo el patrón MVC (modelo-vista-controlador), es decir, de características muy similares a Angular
\item[ \hyperlink{https://www.emberjs.com/}{Ember JS} ] Se trata de un framework web también que representa un cambio de paradigma con respecto a otros frameworks ya que, por ejemplo, para funcionar correctamente nos exige que respetemos totalmente una convenciones de nomenclatura a la hora de nombrar a cada uno de los objetos de la App.
\item[ \hyperlink{https://facebook.github.io/react/}{React} ] Propiedad de Facebook, ReactJS es una librería Javascript de código abierto que ofrece grandes beneficios en performance, modularidad y promueve un flujo muy claro de datos y eventos, facilitando la planeación y desarrollo de apps complejas. 
\end{description}
\end{itemize}


\subsection{Front-end: Boostrap}\label{tecnologias_boostrap}
\hyperlink{https://www.getbootstrap.com/}{Boostrap} es un framework CSS desarrollado inicialmente  por Twitter que permite dar forma a un sitio web mediante librerías CSS que incluyen tipografías, botones, cuadros, menús y otros elementos que pueden ser utilizados en cualquier sitio web.

\imagenflotante{logo-boostrap}{Logo \emph{Boostrap}. Fuente:  \url{getbootstrap.com}.}

\subsection{ Ionic 2}\label{tecnologias_ionic2}
\hyperlink{https://ionicframework.com/}{Ionic2} es un \emph{framework} gratuito y open source para desarrollar aplicaciones híbridas multiplataforma que utiliza HTML5, CSS (generado por SASS) y Cordova como base. Asegura aplicaciones rápidas, con un alto rendimiento y escalables.

\imagenflotante{logo-ionic}{Logo \emph{Ionic}. Fuente:  \url{ionicframework.com}.}

\subsection{Back-end: Node js}\label{tecnologias_nodejs}
\hyperlink{https://nodejs.org/}{Nodejs} es un entorno de programación pensado para realizar funciones de servidor. Permite la construcción de servidores de forma muy sencilla y rápida, además se puede aplicar para otros usos. La ejecución es asíncrona. Esto significa que las funciones no son ejecutadas secuencialmente (si existen dos llamadas consecutivas, no es necesario que acabe la primera llamada para ejecutar la segunda).
Está pensado para ser un gestor de entradas y salidas. No ha sido diseñado para ejecutar gran cantidad de código, sino para realizar comunicaciones muy rápidas y abundantes, tanto con eventos locales como por comunicación por red.
\imagenflotante{logo-nodejs}{Logo \emph{Node JS}. Fuente: \url{nodejs.com}.}

\begin{itemize}
	\item \subsubsection{Alternativa estudiada: PHP}\label{php}
	Un lenguaje ampliamente conocido, con una extensa documentación y multitud de ejemplos. No se integra tan bien como \emph{Node s} con \emph{MongoDB}. 
\end{itemize}


\subsection{Back-end: Express }\label{tecnologias_expressjs}
\hyperlink{http://expressjs.com/}{Express js} es un framework web flexible para \emph{Nodejs} que proporciona un conjunto robusto de características para aplicaciones web y móviles, proporciona una capa delgada de características fundamentales de aplicaciones web. Facilita la creación de API’s gracias la gran variedad de métodos HTTP y middleware que proporciona.
\imagenflotante{logo-express}{Logo \emph{Express}. Fuente:  \url{expressjs.com}.}


\subsection{Base de datos: MongoDB}\label{tecnologias_mongodb}
Para la parte de la base de datos he escogido \hyperlink{https://www.mongodb.com/}{MongoDB}, estamos hablando de un sistema de gestión de bases de datos no relacionales, o \hyperlink{https://es.wikipedia.org/wiki/NoSQL}{noSQL} . Es decir, un base de datos que no tiene tablas, se basa en colecciones de datos.

En las colecciones se almacenan contenidos que pueden tener diferentes campos. Este sistema almacena los datos en documentos de tipo JSON, cosa que puede favorecer la integración con las aplicaciones que trabajen con este tipo de formatos, como es el caso de una API REST. 
\imagenflotante{logo-mongodb}{Logo \emph{MongoDB}. Fuente: \url{mongodb.com}.}


\subsubsection{Detalles: Mongoose}\label{mongoose}
Se ha utilizado un ODM (\emph{Object-document mapping}) para Mongo y Node.js. Nos permite interactuar con la base de datos mediante objetos de JavaScript, y facilita las operaciones CRUD (Create, Read, Update, Delete) y las validaciones de los usuarios.

\subsubsection{Detalles: Mongoskin}\label{mongoskin}
Se ha utilizado un ODM (\emph{Object-document mapping}) para Mongo y Node.js. Nos permite interactuar con la base de datos mediante objetos de JavaScript, y facilita las operaciones CRUD (Create, Read, Update, Delete) y las validaciones de los usuarios.

\begin{itemize}
	\item	\subsubsection{Alternativa estudiada: Cassandra}\label{cassandra}
	\hyperlink{https://www.cassandra.apache.org/}{Apache Cassandra} es una de las mejores opciones cuando la escalabilidad, la alta disponibilidad y la integridad de los datos son las prioridades de un proyecto.
\end{itemize}


\section{Testing}\label{testing}


\subsection{Postman}\label{postman}
Las pruebas básicas se realizaron con  \hyperlink{https://www.getpostman.com}{Postman}. Esta herramienta nos permite construir y gestionar de una forma cómoda nuestras peticiones a servicios \hyperlink{https://es.wikipedia.org/wiki/Transferencia_de_Estado_Representacional}{REST} (Post, Get, etc). Su manejo es realmente intuitivo ya que simplemente tenemos que definir la petición que queremos realizar, es decir introducir la ruta, y pulsar el botón de enviar. Es realmente fácil de usar y está disponible como extensión del navegador Chrome, pero también como aplicación de escritorio.
\imagenflotante{logo-postman}{Logo \emph{Postman}. Fuente:  \url{getpostman.com}.}

\begin{itemize}
	\item	\subsubsection{Alternativa estudiada: SoapUI}
	Mucho más completo que Postman, \hyperlink{https://www.soapui.org/}{SoapUI} tiene una versión gratuita de código abierto y un versión de pago con algunas funcionalidades que hacen que sea mucho más productiva. El motivo de elección del anterior es por que en los cursos realizados aprendí a manejar éste.

\end{itemize}

\subsection{Karma}\label{karma}


\subsection{Jasmine}\label{jasmine}


\section{Documentación}\label{docs}
En este apartado se describen algunas de las herramientas utilizadas para la parte de la documentación del proyecto.

 \subsection{La\TeX}\label{docs_latex}
  LaTex es un sistema de composición de textos, orientado a la creación de documentos escritos que presenten una alta calidad tipográfica. Por sus características y posibilidades, es usado de forma especialmente intensa en la generación de artículos y libros científicos. Como editor para realizar las modificaciones pertinentes he usado \hyperlink{http://www.xm1math.net/texmaker/}{textmaker} .
   \subsubsection{Detalles: texmaker}
   
   Editor de La\TeX para el sistema operativo de MAC. 
  
   \subsubsection{Detalles: ShareLatex}
    \hyperlink{https://es.sharelatex.com/}{ShareLatex}  es un servicio online que nos permite crear y compartir documentos en La\TeX. La característica por la que he empleado este editor es por que  permite la edición conjunta de un proyecto desde el navegador y ofrece diferentes tipos de plantillas según el uso que le vayamos a dar. Por lo tanto para realizar tablas, pruebas de imágenes u otros ensayos antes de copiarlo al documento final las he realizado todas aquí ya que resulta demasiado costoso compilar todo el documento.
  
 \begin{itemize}
 	\item 	\subsubsection{Alternativa estudiada: Open Office}\label{openoffice}
El editor de textos propiedad de Apache fue la primera alternativa elegida, aunque más tarde decidí realizar la memoria en La\TeX. para aprender otra manera de realizar documentos de textos que quizá fuera interesante para el futuro.
\end{itemize}

 \subsection{Mendeley}\label{docs_mendeley}
  \hyperlink{www.mendeley.com/}{Mendeley} es un gestor de referencias bibliográficas. Es una aplicación web y de escritorio, propietaria y gratuita. Permite gestionar y compartir referencias, links, libros o documentos de investigación. Recibí un curso introductorio por parte de la encargada  de gestión de la biblioteca centra de la Universidad de Burgos.
  
  \section{Otras herramientas}\label{otrasherramientas}
En este apartado se describen otras herramientas que se han empleado como son el alojamiento web, el alojamiento de la base de datos, los repositorios u otras tecnologías.

 \subsection{Heroku}\label{otrasherramientas_heroku}
\hyperlink{https://www.heroku.com/}{Heroku}  es un servicio de almacenamiento en la nube que además tiene mecanismos y herramientas para que la puesta en producción de las aplicaciones web sea prácticamente automática.

 \begin{itemize}
 	\item 	 \subsubsection{Alternativas estudiadas: Plataformas en la nube}\label{tnube}
  Plataformas en la nube como Amazon AWS u OpenShift son muy populares hoy en día y permiten prácticamente lo mismo que la aplicación elegida. Se ha elegido Heroku finalmente por la facilidad de la puesta en marcha de la aplicación en producción.
\end{itemize}



  
 \subsection{Mlab}\label{otrasherramientas_mlab}
 \hyperlink{https://mlab.com/}{Mlab} Es una plataforma de base de datos como servicio (\emph{DBaaS}), para alojar y gestionar bases de datos MongoDB. Es gratuita, hasta cierto punto, y muy intuitiva en el uso.


 
  
   \subsection{NPM}\label{otrasherramientas_npm}
Es el gestor de paquetes por defecto para \emph{Nodejs}.
  
   \subsection{Github}\label{otrasherramientas_github}
   Para el control de versiones se ha utilizado \hyperlink{https://www.github.com/}{Github}. 
   
    \subsection{Zenhub}\label{otrasherramientas_zenhub}
    \hyperlink{https://www.zenhub.io/}{Zenhub} es una extensión de Chrome para github. Se utiliza para gestionar proyectos y funciona de manera nativa en la interfaz. Se basa en la metodología ágil yresulta verdaderamente útil a la hora de realizar y gestionar un proyecto, eso sí, hay que cumplir la tareas. 
    
    
\begin{itemize}
 	\item 	 \subsubsection{Alternativa estudiada: Trello}\label{trello}
\hyperlink{https://www.trello.com/}{Trello}.  es un gestor de tareas que permite el trabajo de forma colaborativa mediante tableros compuestos de columnas  que representan distintos estados. Se basa en el método Kanban para gestión de proyectos, con tarjetas que viajan por diferentes listas en función de su estado. Lo comencé utilizando pero al final cambié a Zenhub por su mejor integración con el código. Sin embargo me parece una herramienta muy interesante de la cual hago uso para temas personales.  
\end{itemize}

 \subsection{StackOverFlow}\label{otrasherramientas_dillinger}

\hyperlink{https://stackoverflow.com/}{Stack Over Flowst} es una de las comunidades de desarrolladores más importantes del mundo \cite{articulo2} en la que se responden cuestiones de diferentes lenguajes. Sinceramente es una herramienta fundamental y existe una gran comunidad de usuarios de Angular que crece cada día.


   \subsection{Dillinger}\label{otrasherramientas_dillinger}
    \hyperlink{https://www.zenhub.io/}{Dillinger} es una plataforma \emph{cloud}  basada en HTML5  que permite de manera online modificar código Markdown. Empleada para crear, modificar, añadir mejoras en los archivos \textbf{README.md} .
     
     





\capitulo{5}{Aspectos relevantes del desarrollo del proyecto}

Este apartado pretende recoger los aspectos más interesantes del desarrollo del proyecto: Decisiones que se han tomado, desarrollo, progreso general y problemas que surgieron durante toda la realización del proyecto.


\section{Inicio del proyecto}\label{inicio-proyecto}
Cuando llego la hora de buscar proyecto, entre las  diferentes opciones disponibles estaba \emph{SWAPP}, una herramienta realizada años atrás en la universidad por un alumno que desarrolló una aplicación nativa en Android sobre el mismo tema.

 Desde el primer momento me llamo la atención el proyecto por que resultaba, más allá de ser un proyecto fin de grado, un aplicación que en un futuro podía tener un futuro real. Además se trataba de tecnologías web, no se especificaban cual tras las cuales, un vez desarrollada la parte web, sería posible convertirlas en una aplicación híbrida. Me pareció una idea fascinante, sobre todo por que desconocía plataformas como \emph{Cordova} o \emph{Phohegapp}. Está claro que el principal peso del proyecto se lo lleva la aplicación web, dado que es la parte fundamental y sobre la que se va a sustentar todo lo demás. Por lo que resultaba fundamental escoger un buen framework web, estudiarlo convenientemente y comenzar a realizar un trabajo sólido a partir de ese momemento.
 
  En consecuencia surge el primero de los debates, que tipo de framework usar. La multitud de opciones que existe en el mercado es realmente es extensa, una de las opciones a tener en cuenta para la elección de tu aplicación es comprobar que el framework sigue recibiendo soporte actualmente (Por ejemplo github tiene "escaparates" donde agrupa los frameworks por tipos, algunos repositorios son \hyperlink{ https://github.com/showcases/web-application-frameworks/}{públicos}  y es posible saber la frecuencia con la que se actualizan ). Además una de las mejores cosas del software libre es que puedes interactuar con los propios creadores  o desarrolladores de una tecnologías, abrir \emph{issues} en github y en numerosas ocasiones recibes respuesta. 
  
  
  \section{Formación}\label{formacion}
Desde el primer momento se priorizó la formación por delante de todo ya que el proyecto requiere conocimientos de dos tipos: por una parte conocimientos web en los que el alumno debía especializarse ya que en la Universidad de Burgos no se ve como tal el desarrollo web en ninguna asignatura, por otro lado conocimientos que eran totalmente desconocidos para el alumno como manejo de la parte del servidor, bases de datos NoQL o desarrollo de aplicaciones híbridas. 

Gracias a las nuevas tecnologías tan solo es necesario un ordenador y una buena conexión a internet para tener acceso a conocimiento infinito. Desde hace algunos años esta en auge la formación online y las plataformas web que ofrecen cursos de diversos tipos. Existen numerosísimas plataformas de este estilo, una de las más conocidas es  \hyperlink{www.coursera.org}{Coursera}, que reúne cursos online masivos, abiertos y gratuitos (MOOCS). El curso fue recomendador por mi tutor y sin duda no se equivocaba, desde hace algún tiempo la visibilidad de la parte gratuita se ha reducido pudiendo solo acceder a partes limitadas de los contenidos. Es posible realizar el curso completo pero no obtener feedback de otros usuarios ni tampoco subir a la plataforma los avances de código.

Por lo que en un primer momento se decidió realizar la siguiente especialización compuesta de seis cursos diferentes:

\begin{description}
	\item[\emph{HTML, CSS y JavaScript}] Introducción al desarrollo web con las principales tecnologías. Qué es DOM y CSS.
	\item[ \emph{Front-End Web UI Frameworks and Tools}] Conocimiento de lo que es un framework web: Boostrap. También introducción a los preprocesadores Less, Sass.
	\item[\emph{Front-End JavaScript Frameworks: AngularJS} ]  Introducción y desarrollo con AngularJS.
	\item[ \emph{Multiplatform Mobile App Development with Web Technologies}] Desarrollo UI/UX, el framework Ionic como complemento a AngularJS para  crear aplicaciones hibridas con ayuda de Cordova..
	\item[ \emph{Server-side Development with NodeJS} ] Lado del servidor con Node JS.
	\item[ \emph{Full Stack Web Development Specialization Capstone}] Tabla con detalles de la especialización.
\end{description}


Los cursos pertenecer al curso \hyperlink{https://www.coursera.org/specializations/full-stack}{\emph{full-stack web develompent}} y son bastante completos en lo que a materia se refiere. Si es necesario saber algo de inglés por que sino se hacen bastante pesados.

Después además realicé algún que otro curso en español como son: 

 \section{Frameworks cliente}\label{cliente}
 
 \emph{AngularJs vs Angular 4} está fue una de las preguntas con las que más tiempo perdí al inicio del desarrollo y después de la fase de formación
 
  \section{Tipo de base de datos}\label{base de datos}
  
  Otro de los dilemsa sin duda fue la elección de la base de datos
 
 
\capitulo{6}{Trabajos relacionados}

En esta sección se pretende realizar un acercamiento a los posibles competidores que la empresa podría tener en un potencial mercado, en caso de que llegara a ponerse a la venta. y también a los trabajos en relación ha ésta que hayan existido. En el plan de empresa realizado, del cual se dan más detalles en los anexos, se ha realizado un estudio de la competencia más exhaustivo. 

\section{Trabajos previos}
Indispensable hacer una mención a un trabajo fin de grado previo realizado en este mismo grado y universidad. Aunque no tiene nada que ver el desarrollo, la idea y el tutor sí fueron el mismo luego era necesario nombrarlo.

El trabajo tuvo el nombre de \textbf{ \emph{SWAPP}- Aplicación para gestionar el cambio de turnos y guardias} realizado por el compañero Diegro Prado Nebreda. Desafortunadamente no se encuentra ya en github el proyecto para poder verlo.

\section{Competidores}

\begin{table}[!hbt]
\begin{center}
\begin{tabular}{|l|l|}
\hline
Nombre & \emph{Tipo} \\
\hline
\hyperlink{https://play.google.com/store/apps/details?id=ciesdesign.SaTurnos&hl=es}{Saturnos Pro} & Aplicación Android para smartphones\\
\hline
\hyperlink{https://itunes.apple.com/es/app/cuadraturnos-free-calendario-de-turnos-de-trabajo/id1054129506?mt=8/}{CuadraTurnos Free} & Aplicación Android e IOS para smartphones\\
\hline
\hyperlink{https://myshyft.com/}{Shyft}   & Aplicación Android e IOS para smartphones\\
\hline
\end{tabular}
\caption{Listado de posibles competidores}
\end{center}
\end{table}


En la tabla anterior  se describen tres aplicaciones. Se ha considerado incluir estas tres  por diferentes razones que se expondrán a continuación.

\begin{description}
	\item[\emph{Saturnos pro}] La primera de ellas fue la que más nombro la gente durante las encuestas realizadas para el plan de empresa ya nombrado. Si hay una aplicación usada en este campo en España sin duda es ésta, si bien es cierto que tiene algunas limitaciones como por ejemplo que el único medio para usarla sea un dispositivo móvil con Android. La principal desventaja quizás sea el precio pero también era la más conocida por lo que eso significa que sea usa aunque no masivamente.
	\item[\emph{Cuadraturnos Free}] Fue otra de las aplicaciones nombradas durante la encuesta aunque con mucha menor insistencia, tan solo tres o cuatro personas la conocían. Está disponible en las principales plataformas móviles y resulta fácil de usar pero la interfaz resulta ser un poco antigua .
	\item[\emph{Shyft}] Sin duda es la aplicación más reciente, sólo esta disponible en Estados Unidos pero es el modelo a fijarme como herramienta. La he probado y usado y quizás no resulta tan intuitiva como parece ser.
\end{description}

Nota: Se incorporará un anexo extra con el trabajo realizado en YUZZ sobre la competencia. 

\section{Otros trabajos}
Aunque no directamente relaciones si he encontrado algún trabajo o proyecto que intentaban desarrollar algoritmos para suplir problemas de asignación de turnos. Me a parecido interesante y curioso incluirlos en esta sección.

    \begin{description}
	\item[Trabajo1] Diseño de un modelo de asignación de turnos para la operación
de sistemas de transporte masivo tipo BRT\cite{trabajo1}
	\item[Trabajo2]Algoritmo GRASP \citep{grasp} para resolver el problema de asignación de horarios en empresas de demanda variable \citep{trabajo2}
\end{description}

\capitulo{7}{Conclusiones y Líneas de trabajo futuras}

Todo proyecto debe incluir las conclusiones que se derivan de su desarrollo. Éstas pueden ser de diferente índole, dependiendo de la tipología del proyecto, pero normalmente van a estar presentes un conjunto de conclusiones relacionadas con los resultados del proyecto y un conjunto de conclusiones técnicas. 
Además, resulta muy útil realizar un informe crítico indicando cómo se puede mejorar el proyecto, o cómo se puede continuar trabajando en la línea del proyecto realizado. 



\bibliographystyle{plain}
\bibliography{bibliografia}

\end{document}
